%! Unit 2: Describing One Variable — Summative Assessment — Practice Test --- Study Copy (blank) (Answer Key)
% The practice test is the unit's sample test: tests.mk publishes it to ../sample_test/, and
% shared/unit.mk merges it into the unit student packet.
%
% THE KEY IS GENERATED, NEVER HAND-EDITED:
%   python3 templates/lesson/mkkey.py unit02/tests/practice_test/main.tex \
%       unit02/test_keys/practice_test_key/main.tex unit02/test_keys/practice_test_key/answers.json
% Re-run it after every edit to this file. Part B carries an Answer: slot on every item (the
% quiz-01 model) so the letter lives in the generated key; unit 2 has no unit_cover_key page to
% park the letters on.
\documentclass[10pt]{article}
\usepackage{saar-article}
\usepackage{saar-key}   % pulls in -boxes; do NOT also load -boxes

% Work blocks on a test need handwriting room, not typeset spacing. This moves
% the blank and the key together, so it cannot open a page mismatch.
\setlength{\workrowsep}{5pt}

% Part-divider strip
\newcommand{\parthead}[1]{%
  \vspace{6pt}\noindent
  \begin{headlinebox}{cerulean}{\color{white}\bfseries #1}\end{headlinebox}%
  \vspace{4pt}}

% One horizontal boxplot on a tikz x-axis: \drawbox{y}{min}{Q1}{med}{Q3}{max}{label}
\newcommand{\drawbox}[7]{%
  \draw[charcoal, thick] (#2,#1) -- (#3,#1);
  \draw[charcoal, thick] (#5,#1) -- (#6,#1);
  \draw[charcoal, thick] (#2,#1-0.15) -- (#2,#1+0.15);
  \draw[charcoal, thick] (#6,#1-0.15) -- (#6,#1+0.15);
  \filldraw[fill=frost, draw=cerulean, thick] (#3,#1-0.3) rectangle (#5,#1+0.3);
  \draw[cerulean, very thick] (#4,#1-0.3) -- (#4,#1+0.3);
  \node[left, font=\small\bfseries] at (#2-0.2,#1) {#7};}

\begin{document}

\pageheader{Unit 2: Describing One Variable}{Practice Test --- Study Copy --- Answer Key}
\namedateperiod

\vspace{0.10in}
\begin{remindbox}
\small
\textbf{This is a practice test.} It mirrors the real Unit 2 test in length, format,
and the ideas it covers, but the numbers and contexts differ. Work every problem
with no notes, then check your answers against the key.
\textbf{25 questions, 100 points:} Part A $10$, Part B $12$, Part C $48$, Part D $30$.
\end{remindbox}

% ======================================================================
\parthead{Part A --- Vocabulary \hfill 10 questions, 1 point each}

\small
\textbf{Set 1 --- Displays and Shape.} Write the letter of the matching description
next to each term.

\vspace{2pt}
\begin{minipage}[t]{0.36\linewidth}
\begin{enumerate}[leftmargin=1.9em, itemsep=3.5pt, topsep=1pt]
  \item \ans{C} Distribution
  \item \ans{A} Stemplot
  \item \ans{E} Histogram
  \item \ans{B} Skewed right
  \item \ans{D} Skewed left
\end{enumerate}
\end{minipage}\hfill
\begin{minipage}[t]{0.61\linewidth}
\begin{enumerate}[label=\Alph*., leftmargin=1.9em, itemsep=3.5pt, topsep=1pt]
  \item Each value split into a stem and a leaf, read with a key.
  \item The long tail stretches toward the larger values.
  \item What values a variable takes and how often it takes each one.
  \item The long tail stretches toward the smaller values.
  \item Bars over equal-width bins; each bar's height is a count.
\end{enumerate}
\end{minipage}

\vspace{6pt}
\textbf{Set 2 --- Measures and Outliers.} Same directions.

\vspace{2pt}
\begin{minipage}[t]{0.36\linewidth}
\begin{enumerate}[leftmargin=1.9em, itemsep=3.5pt, topsep=1pt, start=6]
  \item \ans{D} Median
  \item \ans{E} Interquartile range
  \item \ans{B} Standard deviation
  \item \ans{C} Outlier
  \item \ans{A} Resistant
\end{enumerate}
\end{minipage}\hfill
\begin{minipage}[t]{0.61\linewidth}
\begin{enumerate}[label=\Alph*., leftmargin=1.9em, itemsep=3.5pt, topsep=1pt]
  \item Barely moved by one extreme value.
  \item About the typical distance of the values from their mean.
  \item A value below $Q_1 - 1.5 \times \text{IQR}$ or above
        $Q_3 + 1.5 \times \text{IQR}$.
  \item The middle value of the ordered data.
  \item $Q_3 - Q_1$: the spread of the middle half of the data.
\end{enumerate}
\end{minipage}
\normalsize

% ======================================================================
\boxguard[8]
\parthead{Part B --- Multiple Choice \hfill 6 questions, 2 points each}

\small
\begin{enumerate}[leftmargin=2.2em, itemsep=5pt, topsep=2pt, start=11]

\boxguard[9]
\item A histogram of the minutes $60$ Granite Falls students spent at a lunchtime club has
its tallest bars between $0$ and $15$ minutes, then shorter and shorter bars out to $60$
minutes. The distribution is
\begin{enumerate}[label=(\Alph*), leftmargin=1.8em, itemsep=1pt, topsep=1pt]
  \item skewed left, because the tallest bars are on the left.
  \item skewed right, because the long tail stretches toward the larger values.
  \item symmetric, because it has only one peak.
  \item skewed left, because most students spent only a little time.
\end{enumerate}
\hfill \textbf{Answer:} \ans{B}

\boxguard[9]
\item House prices in a town are strongly skewed right by a few very expensive homes. The
best measure of a typical price is
\begin{enumerate}[label=(\Alph*), leftmargin=1.8em, itemsep=1pt, topsep=1pt]
  \item the median, because the few expensive homes barely move it.
  \item the mean, because it uses every price.
  \item the range, because it shows the most expensive home.
  \item the standard deviation, because it measures typical distance.
\end{enumerate}
\hfill \textbf{Answer:} \ans{A}

\boxguard[9]
\item Class P's quiz scores run from $88$ to $96$ with $s = 2.5$. Class Q's run from $60$ to
$95$ with $s = 10$. Which statement is true?
\begin{enumerate}[label=(\Alph*), leftmargin=1.8em, itemsep=1pt, topsep=1pt]
  \item Class P should have the larger $s$, because its scores are higher.
  \item The two standard deviations cannot be compared.
  \item Class Q's scores typically sit farther from their own mean than Class P's.
  \item Class Q scored higher, because its standard deviation is larger.
\end{enumerate}
\hfill \textbf{Answer:} \ans{C}

\boxguard[9]
\item Every value in a data set is increased by $5$. What happens to the mean and the standard
deviation?
\begin{enumerate}[label=(\Alph*), leftmargin=1.8em, itemsep=1pt, topsep=1pt]
  \item Both increase by $5$.
  \item The mean increases by $5$; the standard deviation does not change.
  \item The mean stays the same; the standard deviation increases by $5$.
  \item Neither one changes.
\end{enumerate}
\hfill \textbf{Answer:} \ans{B}

\boxguard[9]
\item One very large outlier is removed from a data set of $15$ values. Which summary
changes the most?
\begin{enumerate}[label=(\Alph*), leftmargin=1.8em, itemsep=1pt, topsep=1pt]
  \item the mean
  \item the median
  \item the first quartile
  \item the interquartile range
\end{enumerate}
\hfill \textbf{Answer:} \ans{A}

\boxguard[9]
\item Two shops each have a median delivery time of $30$ minutes. Shop X's IQR is $4$
minutes and Shop Y's is $12$ minutes. Which is a correct comparison?
\begin{enumerate}[label=(\Alph*), leftmargin=1.8em, itemsep=1pt, topsep=1pt]
  \item The shops perform the same, because their medians are equal.
  \item Shop Y is more consistent, because its IQR is larger.
  \item Shop X is faster, because its IQR is smaller.
  \item Both typically take about $30$ minutes, but Shop X's times are more consistent.
\end{enumerate}
\hfill \textbf{Answer:} \ans{D}

\end{enumerate}
\normalsize

% ======================================================================
\boxguard[8]
\parthead{Part C --- Short Answer \& Computation \hfill 6 questions, 8 points each}

\small
\begin{enumerate}[leftmargin=2.2em, itemsep=9pt, topsep=2pt, start=17]

% ---------------------------------------------------------------- C1 (2.1)
\item Eighteen \textbf{Granite Falls High School} students with part-time jobs recorded the
hours they worked last week.

\vspace{2pt}
\begin{minipage}[t]{0.30\linewidth}
\vspace{0pt}
\renewcommand{\arraystretch}{1.1}
\begin{tabular}{r|l}
0 & 8\;9 \\
1 & 0\;2\;4\;5\;5\;6\;8 \\
2 & 0\;0\;1\;3\;4\;6 \\
3 & 0\;2 \\
4 & 8 \\
\end{tabular}

\vspace{3pt}
\textbf{Key:} $1\,|\,4$ means $14$ hours
\end{minipage}\hfill
\begin{minipage}[t]{0.67\linewidth}
\vspace{0pt}
(a) What does the leaf $8$ on stem $4$ stand for?

\vspace{2pt}
\ans{One student who worked 48 hours}

\vspace{4pt}
(b) Median: \ans{$19$ hours} \quad Range: \ans{$40$ hours}

\vspace{4pt}
(c) Describe the distribution --- shape, unusual features, center, and spread --- in
context.
\end{minipage}

\par\ansline{Roughly symmetric with a slight right tail, a gap from 32 to 48 hours, and a possible outlier at 48.}
\par\ansline{The median is 19 hours, and the hours run from 8 to 48 (range 40 hours).}

% ---------------------------------------------------------------- C3 (2.2)
\boxguard[22]
\item Nine members of the Granite Falls book club read these numbers of books this summer, in
order: $3$, $5$, $6$, $6$, $8$, $9$, $11$, $12$, $21$.

\vspace{2pt}
(a) Find the mean.

\begin{work}
  \bar{x} &= \dfrac{3 + 5 + 6 + 6 + 8 + 9 + 11 + 12 + 21}{9} = \dfrac{81}{9} = 9 \text{ books}
\end{work}

(b) Median: \ans{$8$ books} \quad Mode: \ans{$6$ books}

\vspace{3pt}
(c) Find $Q_1$ and $Q_3$. (Leave the median out of both halves.)

\begin{work}
  Q_1 &= \dfrac{5 + 6}{2} = 5.5 \qquad Q_3 = \dfrac{11 + 12}{2} = 11.5
\end{work}

(d) Five-number summary: \ans{$3$}, \ans{$5.5$}, \ans{$8$}, \ans{$11.5$},
\ans{$21$} \quad IQR: \ans{$6$ books}

% ---------------------------------------------------------------- C4 (2.2)
\boxguard[22]
\item The Granite Falls cafeteria line took these numbers of minutes to clear on five days:
$9$, $11$, $13$, $13$, $14$. The mean is $12$ minutes.

\vspace{2pt}
(a) Complete the table.

\vspace{2pt}
\renewcommand{\arraystretch}{1.3}
\begin{tabularx}{\linewidth}{@{} >{\bfseries}l *{5}{>{\centering\arraybackslash}X} | c @{}}
\rowcolor{frost}
 & & & & & & \textbf{Sum} \\
Value $x$ & $9$ & $11$ & $13$ & $13$ & $14$ & $60$ \\
Deviation $x - \bar{x}$ & \ans{$-3$} & \ans{$-1$} & \ans{$1$} & \ans{$1$} &
  \ans{$2$} & \ans{$0$} \\
Squared $(x - \bar{x})^2$ & \ans{$9$} & \ans{$1$} & \ans{$1$} & \ans{$1$} &
  \ans{$4$} & \ans{$16$} \\
\end{tabularx}

\vspace{2pt}
(b) Find the variance and the standard deviation.

\begin{work}
  s^2 &= \dfrac{16}{5 - 1} = 4 \text{ square minutes} \qquad s = \sqrt{4} = 2 \text{ minutes}
\end{work}

(c) Interpret $s$ in context.

\par\ansline{A typical day's line time is about 2 minutes from the 12-minute mean.}

% ---------------------------------------------------------------- C5 (2.3)
\boxguard[20]
\item The five-number summary of $30$ Granite Falls students' daily phone pickups is
$18$, $42$, $55$, $66$, $110$. The three most extreme values are $18$, $97$, and $110$.

\vspace{2pt}
(a) Find the IQR and both fences.

\begin{work}
  \text{IQR} &= 66 - 42 = 24 \\
  \text{lower fence} &= 42 - 1.5(24) = 6 \qquad \text{upper fence} = 66 + 1.5(24) = 102
\end{work}

(b) Which of $18$, $97$, and $110$ are outliers? \ans{Only $110$}

\vspace{3pt}
(c) In a modified boxplot, where does the upper whisker end? \ans{At $97$}

\vspace{3pt}
(d) The value $18$ looks far from the box. Explain whether it is an outlier.

\par\ansline{No --- 18 is above the lower fence of 6, so it is inside the fences.}

% ---------------------------------------------------------------- C6 (2.3)
\boxguard[26]
\item The graph shows how many sit-ups $100$ Granite Falls students did in one minute.

\vspace{2pt}
\begin{minipage}[t]{0.50\linewidth}
\vspace{0pt}
\begin{tikzpicture}
\begin{axis}[width=7.4cm, height=4.8cm, xmin=10, xmax=70, ymin=0, ymax=100,
  xtick={10,20,...,70}, ytick={0,25,50,75,100}, yticklabel={\pgfmathprintnumber{\tick}\%},
  grid=major, grid style={linegray}, tick label style={font=\scriptsize},
  xlabel={Sit-ups in one minute}, ylabel={Cumulative \%}, label style={font=\scriptsize}]
  \addplot[cerulean, thick, mark=*, mark size=1.3pt]
    coordinates {(10,0) (20,8) (30,25) (40,50) (50,75) (60,92) (70,100)};
\end{axis}
\end{tikzpicture}
\end{minipage}\hfill
\begin{minipage}[t]{0.47\linewidth}
\vspace{2pt}
(a) $Q_1$: \ans{$30$} \quad Median: \ans{$40$}

\vspace{4pt}
$Q_3$: \ans{$50$} \quad IQR: \ans{$20$}

\vspace{4pt}
(b) A student did $60$ sit-ups. What percentile is that? \ans{$92$nd}

\vspace{4pt}
(c) Interpret your answer to (b) in context.
\end{minipage}

\par\ansline{92\% of these 100 students did 60 or fewer sit-ups in one minute.}

% ---------------------------------------------------------------- C8 (2.4)
\boxguard[26]
\item Granite Falls surveyed juniors and seniors about hours of sleep on a school night. The
parallel boxplots share one scale.

\vspace{2pt}
\begin{center}
\begin{tikzpicture}[x=1.5cm, y=1cm]
  \draw[linegray] (3.5,0) -- (10,0);
  \foreach \t in {4,5,...,10} {
    \draw[linegray] (\t,-0.08) -- (\t,0.08);
    \node[below, font=\scriptsize] at (\t,-0.08) {\t};}
  \foreach \t in {4.5,5.5,...,9.5} \draw[linegray] (\t,-0.05) -- (\t,0.05);
  \node[below, font=\scriptsize] at (6.75,-0.45) {Hours of sleep};
  \drawbox{0.7}{4}{5.5}{6.5}{8}{9.5}{Seniors}
  \drawbox{1.6}{5}{6}{7}{7.5}{9}{Juniors}
\end{tikzpicture}
\end{center}

(a) Complete the table.

\vspace{2pt}
\renewcommand{\arraystretch}{1.3}
\begin{tabularx}{\linewidth}{@{} X c c @{}}
\rowcolor{frost}
 & \textbf{Median} & \textbf{IQR} \\
Juniors & \ans{$7$ h} & \ans{$1.5$ h} \\
Seniors & \ans{$6.5$ h} & \ans{$2.5$ h} \\
\end{tabularx}

\vspace{4pt}
(b) Compare the two groups' sleep, with a comparative word for center \emph{and} for spread.

\par\ansline{Juniors' median (7 h) is greater than seniors' (6.5 h), and seniors' sleep is}
\par\ansline{more spread out (IQR 2.5 h vs.\ 1.5 h), so juniors' sleep is more consistent.}

\end{enumerate}
\normalsize

% ======================================================================
\boxguard[10]
\parthead{Part D --- Extended Response \hfill 3 questions, 10 points each}

\small
\begin{enumerate}[leftmargin=2.2em, itemsep=10pt, topsep=2pt, start=23]

% ---------------------------------------------------------------- D1 (2.1, 2.3)
\item Sixty Granite Falls students reported how much they earned last summer. The histogram
and the calculator's 1-Var Stats are below.

\vspace{2pt}
\begin{minipage}[t]{0.56\linewidth}
\vspace{0pt}
\begin{tikzpicture}
\begin{axis}[width=8.2cm, height=4.6cm, ybar interval, xmin=0, xmax=3500, ymin=0, ymax=20,
  xtick={0,1000,2000,3000}, minor xtick={500,1500,2500,3500}, x tick label as interval=false, ytick={0,5,...,20}, ymajorgrids, grid style={linegray},
  tick label style={font=\scriptsize}, xticklabel={\$\pgfmathprintnumber{\tick}},
  scaled x ticks=false, xlabel={Summer earnings (dollars)}, ylabel={Students},
  label style={font=\scriptsize}]
  \addplot[fill=frost, draw=cerulean] coordinates
    {(0,18) (500,16) (1000,10) (1500,7) (2000,4) (2500,3) (3000,2) (3500,2)};
\end{axis}
\end{tikzpicture}
\end{minipage}\hfill
\begin{minipage}[t]{0.40\linewidth}
\vspace{4pt}
\renewcommand{\arraystretch}{1.15}
\begin{tabular}{@{} l r @{}}
\rowcolor{frost}
\multicolumn{2}{@{}l@{}}{\textbf{1-Var Stats}} \\
$\bar{x}$ & $1083$ \\
Sx & $832$ \\
$n$ & $60$ \\
minX & $60$ \\
$Q_1$ & $390$ \\
Med & $860$ \\
$Q_3$ & $1620$ \\
maxX & $3310$ \\
\end{tabular}
\end{minipage}

\vspace{4pt}
(a) Name the shape of the distribution, and justify it from the histogram.

\par\ansline{Skewed right: most students earned under \$1{,}000, and the long tail}
\par\ansline{stretches toward the larger earnings, out to about \$3{,}500.}

(b) The mean (\$1{,}083) is larger than the median (\$860). Explain why, using the shape.

\par\ansline{The mean uses every value, so the few students in the long right tail pull}
\par\ansline{it up; the median uses only the middle position, so it stays lower.}

\boxguard[7]
(c) Which center and which spread should the school report? Give their values in context and
say why.

\par\ansline{The median, \$860, and the IQR, \$1{,}230 (the middle half earned \$390 to \$1{,}620).}
\par\ansline{The data are skewed right, so the tail pulls the mean and SD up.}

(d) A student wrote: \emph{``Skewed left --- most students earned a little, so the peak is on
the left.''} Explain what is wrong.

\par\ansline{The skew is named by the long tail, not the peak. The pile is on the left,}
\par\ansline{but the tail stretches toward larger earnings, so it is skewed right.}

% ---------------------------------------------------------------- D2 (2.2)
\boxguard[30]
\item Two Granite Falls bowlers each rolled five games. \quad
\textbf{Maya:} $180$, $182$, $184$, $186$, $188$ ($\bar{x} = 184$). \quad
\textbf{Jordan:} $110$, $125$, $140$, $150$, $175$ ($\bar{x} = 140$, $s \approx 24.7$).

\vspace{3pt}
(a) Complete the table for Maya, then find her standard deviation.

\vspace{2pt}
\renewcommand{\arraystretch}{1.3}
\begin{tabularx}{\linewidth}{@{} >{\bfseries}l *{5}{>{\centering\arraybackslash}X} | c @{}}
\rowcolor{frost}
 & & & & & & \textbf{Sum} \\
Score $x$ & $180$ & $182$ & $184$ & $186$ & $188$ & $920$ \\
Deviation $x - \bar{x}$ & \ans{$-4$} & \ans{$-2$} & \ans{$0$} & \ans{$2$} &
  \ans{$4$} & \ans{$0$} \\
Squared $(x - \bar{x})^2$ & \ans{$16$} & \ans{$4$} & \ans{$0$} & \ans{$4$} &
  \ans{$16$} & \ans{$40$} \\
\end{tabularx}

\begin{work}
  s^2 &= \dfrac{40}{5 - 1} = 10 \qquad s = \sqrt{10} \approx 3.2 \text{ pins}
\end{work}

(b) A friend says, \emph{``Maya's standard deviation should be the bigger one --- her scores
are way higher.''} Use the formula to explain why the friend is wrong.

\par\ansline{$s$ is built from the deviations --- each score's distance from its own mean.}
\par\ansline{Maya's scores sit within 4 pins of 184; how big the scores are never enters.}

(c) Next week every one of Maya's five scores is exactly $10$ pins higher.
New mean: \ans{$194$} \quad New standard deviation: \ans{$3.2$}

\vspace{3pt}
Explain why the standard deviation came out that way.

\par\ansline{Adding 10 to every score adds 10 to the mean, so not one deviation changes.}

(d) Which bowler is more consistent? Justify with numbers.

\par\ansline{Maya: $s \approx 3.2$ pins against Jordan's 24.7, so her games stay close together.}

% ---------------------------------------------------------------- D3 (2.4)
\boxguard[30]
\item Granite Falls's science department tested $12$ batteries of each of two brands. Here
are the lifetimes, in hours, from the calculator.

\vspace{3pt}
\renewcommand{\arraystretch}{1.25}
\begin{tabularx}{\linewidth}{@{} X *{7}{c} @{}}
\rowcolor{frost}
\textbf{Brand} & $\bar{x}$ & \textbf{Sx} & \textbf{minX} & $Q_1$ & \textbf{Med} & $Q_3$ &
  \textbf{maxX} \\
\midrule
Brand A & $40$ & $2.1$ & $36$ & $38.5$ & $40$ & $41.5$ & $44$ \\
Brand B & $40$ & $7.8$ & $26$ & $34$ & $40$ & $46$ & $52$ \\
\end{tabularx}

\vspace{4pt}
Both distributions are roughly symmetric.

\vspace{3pt}
(a) Check Brand B for outliers with the $1.5 \times \text{IQR}$ rule.

\begin{work}
  \text{IQR} &= 46 - 34 = 12 \\
  \text{fences} &= 34 - 18 = 16 \quad\text{and}\quad 46 + 18 = 64
\end{work}

\par\ansline{Every value lies between the fences (26 to 52), so Brand B has no outliers.}

(b) The department head says, \emph{``Same median, same mean --- it doesn't matter which we
buy.''} Respond.

\par\ansline{When the centers match, the spread decides. Brand B's lifetimes are much}
\par\ansline{more spread out, so some B batteries die far sooner than any A battery.}

(c) Write a comparison of the two brands --- shape, outliers, center, and spread --- using
comparative words.

\par\ansline{Both brands are roughly symmetric with no outliers, and both centers are 40 hours.}
\par\ansline{Brand B is much more spread out (Sx 7.8 h, range 26 h) than Brand A}
\par\ansline{(Sx 2.1 h, range 8 h), so Brand A's lifetimes are more consistent.}

(d) A flashlight must run at least $35$ hours. Which brand should the department buy?
Justify with the table.

\par\ansline{Brand A. Its minimum is 36 hours, so every A battery lasted at least 35 hours;}
\par\ansline{Brand B's $Q_1$ is 34, so at least a quarter of B batteries fall short.}

\end{enumerate}
\normalsize

\end{document}
